\documentclass[anonymous,notoc]{gmcm2026}
\usepackage{placeins,seqsplit,calc}
\providecommand{\tightlist}{\setlength{\itemsep}{0pt}\setlength{\parskip}{0pt}}
\providecommand{\pandocbounded}[1]{#1}
\setlength{\emergencystretch}{2em}
\makeatletter\setlength{\@fptop}{0pt}\makeatother
\gmcmsetup{title={面向多核 NPU 的语义一致结构化切分与流水调度},year=2026,edition={第二十三届},problem=A}
\begin{document}
\maketitle
\begin{abstract}
多核神经网络处理器的计算图切分同时改变并行度、数据搬运、存储驻留和流水线就绪关系。切分后看似均衡的计算负载，仍可能被子图屏障、共享带宽或缓存占用抵消。针对赛题给定的三种硬件语义，本文以统一的结构化求解框架为主线，将原始计算图、可提交方案、实际执行事件和评价指标分层建模，在统一合法性约束下构造少量具有明确结构依据的候选，再用与官方结果对应的评价闭环确定输出。

针对问题一，建立包含子图依赖和同核执行顺序的增广任务图，以屏障代价约束切分粒度。算法依据连通结构、重组件及分叉前沿构造候选，并在满足波次秩条件的区域中采用``独立双支路---汇合尾部''三片援助，增加并行机会而不引入依赖环。对问题二，将同核子图合并执行带来的复用写成按触及 Core 数计费的超图模型，结合链与汇合感知的空隙插入、带锚点的两核局部最小割及负载守卫细化，最多评价三个去重完整方案。对问题三，利用私有归约树的结构，证明按子树内部峰值与保留结果之差排序，在非交错的局部存储模型中最小化内部存活前沿；将该次序作为受限候选，并以必要执行下界剪枝、官方 Makespan 严格改进准则接受。

在全部 100 个官方计算图、每图 1～5 核的固定配置下，三种固定求解器均得到 500 个合法结果。问题一和问题二在 2～5 核上的平均加速比分别为 1.9516、2.7514、3.4786、4.0553 和 2.3167、3.2356、3.9711、4.5498。问题三在 5 核上的平均基线加速比为 4.7576；对相同计划进行无 L2 与只读 Cache 配对评价，逐图 Cache 加速比的均值为 1.0829，按字节计算的命中率均值为 31.36\%。问题一 5 核求解墙钟均值为 3.522 s，问题二全部 500 格均值为 4.029 s；二者均包含在线构造与选优，属于共享主机上的实测结果。

研究表明，依赖结构、数据复用和物理执行次序需要联合处理，单独减少搬运量或提高 Cache 命中率不能保证完成时间缩短。本文以实际计算图与方案的对应关系解释切分机制，以核心和 Pipe 时间线检验重叠效果，并保留固定构造的负结果及证明适用边界。所给算法提供了可复现的结构化求解方法，不依赖逐用例历史赢家拼接，也不将局部模型最优性外推为一般多核调度的全局最优性。
\keywords{多核 NPU；计算图切分；流水调度；超图复用；归约森林；语义一致性}
\end{abstract}
\section{问题重述与分析}\label{ux95eeux9898ux91cdux8ff0ux4e0eux5206ux6790}

\subsection{问题背景}\label{ux95eeux9898ux80ccux666f}

神经网络计算图通常同时含有矩阵计算、向量运算和数据搬运。一个 NPU 核心内的 Cube、Vector、搬入和搬出单元可在满足依赖时并行工作；多个核心又共享核外 DDR。因此，计算图上的一条边不仅表示逻辑先后关系，还可能在切分后转化为存储分配、数据复制与跨核同步。图切分的效果需要在最终执行过程中判断。

传统多处理器调度中的工作量平衡和关键路径分析提供了基础工具，但资源增加或局部次序变化可能引起调度反常现象{[}2{]}。本题进一步包含有限的 L1／UB、固定核内调度规则、共享 DDR 和只读 FIFO Cache。一个子图的执行时间会随其所处计划和其他核心的活动改变，因而不能简单地给每个子图赋一个固定标量权重，再把问题化为普通的独立任务装箱。

本文遵循赛题规定{[}1{]}，以未修改评价程序在固定配置下计算的总完成时间 Makespan 为主要质量指标，同时报告额外 DDR 搬运量，问题三另报告按字节统计的 Cache 命中率。生成方案的程序墙钟是另一类指标。模拟周期衡量方案，墙钟衡量求解程序，两者分别评价。

\subsection{输入、输出及三问的关系}\label{ux8f93ux5165ux8f93ux51faux53caux4e09ux95eeux7684ux5173ux7cfb}

输入是 Op--Tensor 二部有向无环图。张量结点携带位置及字节数，操作结点携带类型、所属 Pipe 及周期数；依赖边连接张量和操作。输出由两个字段构成：\texttt{node\_to\_subgraph} 将全部非 COPY 操作映射到子图编号；\texttt{core\_schedules} 为每个 Core 列出子图的执行顺序。算法可以决定切分、分核及子图顺序，不能直接修改官方核内调度程序、原始计算量、缓存容量或带宽。

三问共享这个输出接口，但对方案的解释不同。

{\def\LTcaptype{table} % do not increment counter
\begin{longtable}[]{@{}
  >{\raggedright\arraybackslash}p{(\linewidth - 6\tabcolsep) * \real{0.2500}}
  >{\raggedright\arraybackslash}p{(\linewidth - 6\tabcolsep) * \real{0.2500}}
  >{\raggedright\arraybackslash}p{(\linewidth - 6\tabcolsep) * \real{0.2500}}
  >{\raggedright\arraybackslash}p{(\linewidth - 6\tabcolsep) * \real{0.2500}}@{}}
\toprule\noalign{}
\begin{minipage}[b]{\linewidth}\raggedright
比较对象
\end{minipage} & \begin{minipage}[b]{\linewidth}\raggedright
问题一
\end{minipage} & \begin{minipage}[b]{\linewidth}\raggedright
问题二
\end{minipage} & \begin{minipage}[b]{\linewidth}\raggedright
问题三
\end{minipage} \\
\midrule\noalign{}
\endhead
\bottomrule\noalign{}
\endlastfoot
执行 Task & 一个子图一个 Task & 同核全部子图合并为一个 Task & 同问题二 \\
同核跨子图数据 & 经过 DDR，Task 间不保留 & 可在核内驻留复用 & 同问题二 \\
跨核同步 & 前驱 Task 完成后等待 & 对应 COPY\_OUT 完成后等待 & 同问题二，命中也不取消依赖 \\
核外只读 Cache & 无 & 无 & 所有核心共享，FIFO 淘汰 \\
主要新增决策影响 & 切分粒度与屏障 & 复用、容量和发射次序 & 查询次序、命中与双带宽池竞争 \\
\end{longtable}
}

问题一要求降低切分造成的启动、切换和边界复制代价，同时保持足够并行度。问题二放开同核驻留能力，原来每个子图独立的容量压力转变为跨子图的存活区间压力，分核与排序必须联合考虑。问题三进一步允许跨核读共享，其收益取决于请求是否在正确时间到达、数据是否仍在 Cache 以及其他请求的竞争状态。本文将这一递进关系贯穿模型、算法和图解，而不是分别构建三个互不关联的评分公式。

\begin{figure}[!htbp]
\centering
\includegraphics[width=\linewidth,height=.84\textheight,keepaspectratio,trim=0 185 0 0,clip]{figures/fig-progression-zh.png}
\caption{三问执行语义的递进。相同图片段和分核方式下，P1 的独立 Task 边界需要 COPY；P2 允许核内驻留；P3 为可查询的 COPY\_IN 增加共享只读 Cache。}
\label{fig:progression}
\end{figure}

\subsection{难点分解与研究路径}\label{ux96beux70b9ux5206ux89e3ux4e0eux7814ux7a76ux8defux5f84}

\textbf{第一，逻辑图合法不等于分核次序合法。} 把操作收缩为子图后，需要重新检查子图间有向环；加入各核的执行顺序后，还可能形成新的等待环。合法性必须在完整提交方案上判定，不能只看每一块内部连通或每一条原始依赖。

\textbf{第二，边界代价具有复用和时序效应。} 同一个输入被多个操作消费时，数据搬运次数与消费者触及的 Core 集合有关，并不总是与切断的操作依赖边数成正比。与此同时，保留数据占用缓存，消除一次复制也可能延长存活区间，挤出其他数据。

\textbf{第三，真实流水重叠由动态资源共同决定。} 计算单元空闲可能源于数据尚未到达、跨核释放未到、存储申请次序未轮到或其他执行前驱。仅凭时间线上的白色间隙不能断言其原因；需要把图、编译事件和 trace 关联起来。

\textbf{第四，复杂模型必须与求解成本相匹配。} 题面建议避免暴力迭代搜索，以单个用例 5～10 min 给出结果为佳。本文采用结构识别、局部精确子问题与有界候选组合，使每一次评价都有明确用途；不把``限定搜索次数''本身作为方法合理性的论据。

由此形成全文的研究路径：先明确可提交自由度与执行语义；再提取适合构造候选的结构量；用合法性和有适用域的必要下界排除无效候选；最后在完整方案上评价、选择并保留可复核结果。三问共用该逻辑，分别引入屏障、复用和缓存次序三个层次的机制。

\FloatBarrier

\section{模型假设与符号说明}\label{ux6a21ux578bux5047ux8bbeux4e0eux7b26ux53f7ux8bf4ux660e}

\subsection{依据题面采用的条件}\label{ux4f9dux636eux9898ux9762ux91c7ux7528ux7684ux6761ux4ef6}

\begin{enumerate}
\def\labelenumi{\arabic{enumi}.}
\tightlist
\item
  原始图及官方配置保持不变。张量和操作 ID 在同一图内唯一，字节数及周期字段非负，原始依赖图无环；不删除计算、不改变操作数值语义。
\item
  核心同构，每核 M、V、MTE2、MTE3 各有单槽执行资源。同核不同 Pipe 与不同 Core 可并发，但发射仍须满足依赖、Task 激活和存储申请条件。
\item
  核内 L1 与 UB 容量分别为 524288 B 和 131072 B，两个存储池独立检查。所有 DDR 读写共享 60 B/cycle，问题三 Cache 容量为 1048576 B，命中读共享独立的 250 B/cycle 带宽池。
\item
  问题一的同核 Task 切换等待为 100 cycles，跨核前驱 Task 等待为 1000 cycles；问题二、三跨核 COPY 同步等待为 500 cycles。具体等待作用于题面规定的完成与释放事件。
\item
  官方核内调度、换入换出、存储申请顺序、整数周期及事件处理规则均属于执行语义。最终结果以原评价程序为准，不以理想并行模型替换。
\end{enumerate}

以上是题目条件的建模化表达，不额外假设带宽无限、Cache 完美、计算与搬运必然完全重叠或所有任务耗时与环境无关。后文树排序、超图计费和候选下界所需的附加条件均在命题处单独给出；结构识别不通过时不使用相应结论。

\subsection{对象层次与术语约定}\label{ux5bf9ux8c61ux5c42ux6b21ux4e0eux672fux8bedux7ea6ux5b9a}

为避免同一个``图''字承载不同约束，全文区分下列对象。

{\def\LTcaptype{table} % do not increment counter
\begin{longtable}[]{@{}
  >{\raggedright\arraybackslash}p{(\linewidth - 6\tabcolsep) * \real{0.2500}}
  >{\raggedright\arraybackslash}p{(\linewidth - 6\tabcolsep) * \real{0.2500}}
  >{\raggedright\arraybackslash}p{(\linewidth - 6\tabcolsep) * \real{0.2500}}
  >{\raggedright\arraybackslash}p{(\linewidth - 6\tabcolsep) * \real{0.2500}}@{}}
\toprule\noalign{}
\begin{minipage}[b]{\linewidth}\raggedright
层次
\end{minipage} & \begin{minipage}[b]{\linewidth}\raggedright
对象
\end{minipage} & \begin{minipage}[b]{\linewidth}\raggedright
保留的信息
\end{minipage} & \begin{minipage}[b]{\linewidth}\raggedright
不能据此直接声称
\end{minipage} \\
\midrule\noalign{}
\endhead
\bottomrule\noalign{}
\endlastfoot
输入层 & 原始二部图 & 操作、张量、大小、位置及原依赖 & 已经有执行次序 \\
计算投影层 & 计算依赖 DAG & 保守保留的原计算因果 & 完整 COPY/MEM 关系 \\
决策层 & 子图、分核及核序 & 参赛者能提交的两字段方案 & 直接控制操作发射时刻 \\
结构层 & 增广子图 DAG & 数据依赖与全部核链 & 容量及动态执行成功 \\
编译层 & Task 与扩展事件图 & COPY、spill、FIFO、存储申请规则 & 所有事件具有环境无关的时长 \\
运行层 & 事件轨迹 & 发射、完成、驻留、共享服务与 Cache & 单一聚合分数足以恢复整个轨迹 \\
观测层 & M、DDR、命中率及墙钟 & 各自明确定义的结果 & 不同指标可以相互替代 \\
\end{longtable}
}

``子图''是方案中的划分对象；``Task''是评价器按场景构造的执行对象，二者只在问题一中一一对应。``原始张量''\,``编译后物理存储实例''和``Cache 查询使用的逻辑标识''分别按所在层解释，不能仅因显示名称相同就视作同一内存对象。

用 \(\Theta=(G,q,K,C,H,R,I)\) 标记一次评价的语义身份：\(C\) 为配置，\(H\) 为官方源码身份，\(R\) 为相关运行环境，\(I\) 为入口、解析及调用契约。性能环境另记录硬件和并发。固定源码不能单独证明跨解释器或浮点运行时完全一致。

形式合法的输入集合、增广无环的结构方案集合和官方成功执行的方案集合分开记为 \(\mathcal P_{\mathrm{in}}\)、\(\mathcal P_{\mathrm{ord}}\)、\(\mathcal P_{\mathrm{exec}}\)。Makespan 只在成功结果上定义。官方拒绝、外部超时、取消、服务错误和未评价均保留各自状态；超时并不证明方案必然无效，也不填成零周期或无穷成绩参加均值。

本文的``精确''必须带对象：可以是某个费用公式精确、固定局部子问题精确、特定输入域的评分等价，或完整执行结果逐字段一致。四者不能省略限定后统一称为``精确求解''。类似地，``下界''必须说明约束全问题、受限方案族、一个固定方案还是某个理想模型。

\subsection{主要符号}\label{ux4e3bux8981ux7b26ux53f7}

{\def\LTcaptype{table} % do not increment counter
\begin{longtable}[]{@{}lll@{}}
\toprule\noalign{}
符号 & 含义 & 单位或范围 \\
\midrule\noalign{}
\endhead
\bottomrule\noalign{}
\endlastfoot
\(G=(O\cup T,E)\) & 原始操作---张量二部图 & DAG \\
\(V\subset O\) & 非 COPY 原始操作集合 & 待划分对象 \\
\(b_t,\ell_t\) & 张量 \(t\) 的大小及位置 & B；DDR/L1/UB \\
\(d_v,p_v\) & 操作计算时长及 Pipe & cycles；M/V \\
\(K\) & Core 数 & \(1,\ldots,5\) \\
\(\pi(v)\) & 操作所属子图 & 非负整数 \\
\(a(s)\) & 子图所属核心 & \(0,\ldots,K-1\) \\
\(\sigma_k\) & 第 \(k\) 核的子图序列 & 有序列表 \\
\(P=(\pi,a,\sigma)\) & 完整可提交方案 & 两字段 JSON 的数学表示 \\
\(\mathcal C_q(G,P)\) & 场景 \(q\) 的 Task 与事件编译 & 含资源和同步约束 \\
\(M_q(P)\) & 官方 Makespan & cycles \\
\(D_q(P)\) & 官方额外 DDR 搬运量 & B \\
\(B_i\) & 第 \(i\) 图官方单核启发式基线 & cycles \\
\(S_{q,i,K}\) & 第 \(i\) 图的基线加速比 & \(B_i/M_q(P_{i,K})\) \\
\(G_{i,K}\) & P3 同计划 Cache 加速比 & \(M_2(P_{i,K})/M_3(P_{i,K})\) \\
\(H_{i,K}\) & P3 按字节计算的 Cache 命中率 & \([0,1]\) \\
\(L,U\) & 同一优化域的必要下界与可行上界 & cycles \\
\(T_{\mathrm{solve}}\) & 读入至方案落盘的求解墙钟 & s \\
\end{longtable}
}

后文出现 \(d_v=\max(1,\mathrm{cycles}_v)\) 时，专指与官方计算投影一致的正整数执行时长；输入的原始字段仍允许为零。图中张量用圆角小框，操作用矩形，子图用带编号外框，核心用水平泳道。该约定用于区分``数据是什么''\,``计算做什么''\,``谁先谁后''和``放在哪个 Core''四种信息。

\subsection{指标定义与比较边界}\label{ux6307ux6807ux5b9aux4e49ux4e0eux6bd4ux8f83ux8fb9ux754c}

对每个核数，平均加速比定义为

\[\bar S_{q,K}=\frac1{100}\sum_{i=1}^{100}\frac{B_i}{M_q(P_{i,K})}.\]

问题一、二的题目主曲线在 \(K=1\) 处按单核基线定义为 1；算法输出的优化单核方案可以另报 \(B_i/M_q(P_{i,1})\)，不得把二者混写。由于 \(B_i\) 是启发式基线而非单核最优值，加速比大于 Core 数不自动构成异常。

问题三的 Cache 比较使用同一个 \(P_{i,K}\)、图与配置，分别启用和关闭 L2。正文同时给出 \(\operatorname{mean}(B/M_2)\)、\(\operatorname{mean}(B/M_3)\) 与 \(\operatorname{mean}(M_2/M_3)\)；前两条均值曲线之商通常不等于第三个量。

额外搬运量采用官方结果中的``调度后 COPY 字节数减去原始图 COPY 字节数''。问题三的该字段仍沿用官方 COPY 计费口径，不因命中而任意改计为零；因此不能把它直接称为 Cache 之后的实际物理 DDR 流量。Cache 命中率使用命中字节除以查询字节；跨用例给出的均值为逐图命中率的算术平均，不是把所有图字节拼接后得到的流量加权命中率。

求解墙钟包含输入读取、预处理、候选生成、在线评分、回退和方案落盘；在外部完成的最终官方复评单独报告。在线使用官方评价的时间已经属于求解墙钟，不能再次从中剔除。共享主机的并发批次墙钟、逐格进程时间和独占机器延迟是不同量，本文不相互替代。

\FloatBarrier

\section{统一模型与本文研究框架}\label{ux7edfux4e00ux6a21ux578bux4e0eux672cux6587ux7814ux7a76ux6846ux67b6}

\subsection{从计算图到执行事件的分层模型}\label{ux4eceux8ba1ux7b97ux56feux5230ux6267ux884cux4e8bux4ef6ux7684ux5206ux5c42ux6a21ux578b}

设 \(V\) 为全部非 COPY 原始操作。切分映射 \(\pi:V\to\mathcal S\) 的非空原像形成子图集合，分核函数 \(a:\mathcal S\to\{0,\ldots,K-1\}\) 与序列 \(\sigma_k\) 满足：每个子图恰出现在一个序列中，且恰出现一次。由原始张量依赖导出子图关系 \(E_{\mathcal S}\)，再加入同核序列中相邻子图之间的顺序边，得到

\[D^+(P)=(\mathcal S,E_{\mathcal S}\cup E_\sigma).\]

提交方案至少须满足覆盖、唯一调度与增广图无环。合法方案在场景 \(q\) 下经官方规则生成 Task、边界 COPY、换入换出和同步事件，再由事件模拟器得到结果。将该过程写为

\[P\xrightarrow{\mathcal C_q}\mathcal E_q\xrightarrow{\mathcal R_q}(M_q,D_q,H_q).\]

这里 \(\mathcal E_q\) 不只是一个固定权重的 DAG：共享 DDR 中操作的剩余工作量会随并发请求数改变，Cache 状态又由发射、完成与淘汰事件更新。把它写成``编译---执行''两层，是为了区分静态依赖和动态资源竞争，并未假设可以在编译时精确确定全部事件时刻。

优化目标是

\[\min_{P\in\mathcal F_q(G,K)} M_q(P),\qquad\text{并报告 }D_q(P)\text{ 及问题三的 }H_q(P).\]

其中 \(\mathcal F_q\) 还受真实容量和事件合法性约束。本文部分选择器用 \((M,D)\) 字典序决策：优先降低 Makespan，时间相同再比较搬运量。这是使候选选择确定化的算法规则，不是题面给出的加权评分公式。P3 Forest 采用更严格的时间改善接受条件，时间持平时保留当前方案。

\subsection{以一致语义连接模型、构造和验证}\label{ux4ee5ux4e00ux81f4ux8bedux4e49ux8fdeux63a5ux6a21ux578bux6784ux9020ux548cux9a8cux8bc1}

本文围绕三问逐步扩展的执行语义建立统一研究框架：\textbf{在同一个输入和执行契约下，选择能够保留当前决策所需信息的结构表示，并明确哪些简化只支持构造、哪些结论足以支持拒绝、哪些收益仍须完整评价。} 它落实为三项要求。

第一，表示的精细程度随决策而变化。判断分叉独立性使用计算依赖图；估计复用使用张量超边及触及 Core 集合；安排归约顺序使用子树内部峰值和根结果大小。每种表示都有使用域，不能把一个只描述字节数的状态当作足以预测未来 FIFO、Cache 和 DDR 行为的完整机器状态。

第二，机制可以并行组合。问题一的组件装箱与分叉援助、问题二的基线与 gap 候选、问题三的现有构造与森林排序，都可以从同一输入产生竞争方案。总体框架因此是具有守卫与回退的机制 DAG，而不是所有样本都必须经过同一串变换的长链。

第三，候选构造与接受证据分开。启发式可以提出可能更好的方案；合法性证书只能说明不会因相应约束失效；必要下界可以在限定域中拒绝不可能胜出的候选；最终质量由完整执行确定。这样可以同时使用简洁的结构模型和忠于物理语义的评价，不要求每一步都重建整个模拟器。

\begin{figure}[!htbp]
\centering
\includegraphics[width=\linewidth,height=.84\textheight,keepaspectratio,trim=0 100 0 0,clip]{figures/fig-framework-zh.png}
\caption{本文研究框架。三问各自根据结构提出候选，共同通过合法性、适用的必要界与完整执行评价决定接受；各机制按适用条件启用。}
\label{fig:framework}
\end{figure}

\subsection{增广无环与局部重构的安全边界}\label{ux589eux5e7fux65e0ux73afux4e0eux5c40ux90e8ux91cdux6784ux7684ux5b89ux5168ux8fb9ux754c}

\textbf{命题 3.1（秩函数证书）。} 若存在整数函数 \(r:\mathcal S\to\mathbb Z\)，使增广任务图每条边 \(s\to t\) 都满足 \(r(s)<r(t)\)，则 \(D^+(P)\) 无有向环。

\textbf{证明。} 若有环 \(s_1\to\cdots\to s_m\to s_1\)，沿环严格递增给出 \(r(s_1)<r(s_1)\)，矛盾。\(\square\)

该命题说明局部构造可通过保持外部秩、为内部新子图细化秩来证明安全；它不说明方案具有容量可行性或更短执行时间。实际算法仍进行完整图检查。

一个容易误用的替代条件是``每个子图路径闭合''。即使两个子图内部各自没有绕出再返回的路径，也可能存在两条原始依赖 \(a_1\to b_1\) 和 \(b_2\to a_2\)；将 \(a_1,a_2\) 分在一组、\(b_1,b_2\) 分在另一组后，商图出现双向边。原图无环而商图有环，因此局部闭合检查不能替代完整增广无环检查。此反例也解释了为什么切图可视化必须同时展示跨边界依赖，而不能只画互不相交的彩色框。

\subsection{必要下界、可行上界与局部最优性}\label{ux5fc5ux8981ux4e0bux754cux53efux884cux4e0aux754cux4e0eux5c40ux90e8ux6700ux4f18ux6027}

对保留全部原始计算、满足每核每 Pipe 单槽条件的执行，令

\[W_M=\sum_{v:p_v=M}d_v,\qquad W_V=\sum_{v:p_v=V}d_v.\]

在确认原始计算依赖被真实执行保留的投影图上，设 \(L_{\mathrm{cp}}\) 为最长计算路径，则有基础必要下界

\[L_{\mathrm{basic}}=\max\left\{L_{\mathrm{cp}},\left\lceil\frac{W_M}{K}\right\rceil,\left\lceil\frac{W_V}{K}\right\rceil,\max_vd_v\right\}\le M.\]

工作量项分别针对 M、V 计算。两条 Pipe 可并行，因此不能将 \(W_M+W_V\) 作为同一服务器工作量直接除以 \(K\)。投影也不能任意收缩经过 COPY 的路径，再把由此得到的关系解释为原始计算的直接边。

对于固定候选 \(P\)，若能额外确定同核同 Pipe FIFO 和必要跨核释放关系，可得到更强的候选下界 \(L(P)\)。当当前成功方案的完成时间为 \(U\)，且 \(L(P)\ge U\)，即可拒绝严格改善不可能的这个候选。这里 \(L(P)\) 约束的是已固定的候选，通常不是所有方案最优值的下界。把候选剪枝界拿来评估全问题最优差距，会混淆优化域。

若 \(L\le\mathrm{OPT}\le U\) 确实来自同一优化域，则

\[\frac{U}{\mathrm{OPT}}-1\le\frac UL-1.\]

这给出保守的质量保证。\(U/L-1\) 小可以证明距离最优较近，值大则可能来自算法不足，也可能来自下界放松，不能被解释为一定可实现的改进空间。

\subsection{状态压缩为什么需要更强的证明}\label{ux72b6ux6001ux538bux7f29ux4e3aux4ec0ux4e48ux9700ux8981ux66f4ux5f3aux7684ux8bc1ux660e}

上述统一建模进一步提出一个问题：能否用比完整方案更小的状态，正确决定后续修改及其结果？这一问题有助于限定编译复用和局部响应模型的解释范围，但当前三种主算法并未实现一般的最小未来充分状态。

先固定允许的动作集合 \(\mathcal A\) 及每个动作的参数语义，例如在一个声明的结构域内切分、合并或移动子图；动作 \(a\) 的定义域为 \(D_a\)。选定在所声明研究域上有定义的观测 \(\Phi\)：在成功执行域可以只含 Makespan；若域含未成功状态，则另用明确的结果状态观察，不能把未定义的 M 填为数值。不同输出投影对应不同等价关系。令 \(L(P)\) 是从 \(P\) 出发可执行的有限动作词集合，包含空字。未来行为等价要求

\[L(P)=L(Q),\qquad \Phi(w(P))=\Phi(w(Q)),\quad\forall w\in L(P).\]

仅比较双方``碰巧都合法''的动作，会漏掉可用动作不同的情况；仅比较当前 \(M(P)=M(Q)\)，也不能排除下一次切分后 FIFO 或存储行为不同。因此当前分数相同并不自动产生可以安全搜索的商空间。

\textbf{命题 3.2（受限域的未来充分表示）。} 给定压缩 \(S\)，若对每个动作存在可用性判定 \(d_a\)、状态更新 \(\bar a\) 及当前观测函数 \(o\)，并且对声明域内全部方案满足

\[P\in D_a\iff d_a(S(P)),\quad S(a(P))=\bar a(S(P)),\quad\Phi(P)=o(S(P)),\]

则 \(S(P)=S(Q)\) 蕴含二者未来行为等价。

\textbf{证明。} 长度为零的词由观测条件成立。假定任意长度 \(j\) 的共同可执行词作用后压缩状态仍相等，则可用性条件保证下一动作在两侧同时可用或同时不可用；更新条件使执行后的压缩状态相等，再由观测条件得到输出一致。对词长归纳，合法语言与全部未来观测相同。\(\square\)

该命题给出的是充分条件。要称为``最小''，还需证明不同压缩状态能够被某个未来词区分，或证明所得划分最粗；把完整方案编码进一个整数并未减少信息成本。本文对实际优化只使用已经核定的较弱对象：完全相同计划去重、满足键条件的局部编译复用、受限模型的响应表示和局部候选构造。数学深度体现在区分这些保证并证明需要的条件，不依靠把所有复用统称为低维商。

\subsection{有界候选与求解复杂度的组织}\label{ux6709ux754cux5019ux9009ux4e0eux6c42ux89e3ux590dux6742ux5ea6ux7684ux7ec4ux7ec7}

统一求解流程为：解析与结构识别；构造基准方案；按输入结构产生有限候选；去重并检查合法性；使用适用的必要条件剪枝；评价剩余完整方案；输出最佳已确认方案。算法没有用例编号到赢家的查表分支，历史实验只用于设计固定规则。

设解析和结构构造总成本为 \(C(G,K)\)，完整图合法性检查为 \(O(|V|+|E_V|+|\mathcal S|)\)，实际评价次数为 \(b\)，各次评价成本为 \(T_j\)，则总体成本可记为

\[T_{\mathrm{solve}}=O\!\left(C(G,K)+b(|V|+|E_V|+|\mathcal S|)\right)+\sum_{j=1}^{b}T_j+T_{\mathrm{IO}}.\]

该表达式不把评价器时间隐藏在一个常数中。本文同时给出候选数量上限、实际调用账和端到端墙钟；对于内部精确最小割等子过程，再在对应章节说明其规模及调用上界。有限候选保证评价次数有界，求解高效还依赖候选规模、结构识别成本及资源上限。

\FloatBarrier

\section{问题一：Task 屏障感知的结构化切分}\label{ux95eeux9898ux4e00task-ux5c4fux969cux611fux77e5ux7684ux7ed3ux6784ux5316ux5207ux5206}

\subsection{从负载均衡转向边界与并行度的联合选择}\label{ux4eceux8d1fux8f7dux5747ux8861ux8f6cux5411ux8fb9ux754cux4e0eux5e76ux884cux5ea6ux7684ux8054ux5408ux9009ux62e9}

问题一中，每个子图成为独立 Task。若将一个大子图拆成两个小子图，可能把工作交给空闲核心，也会产生边界数据搬运和新的激活条件。即使两个子图位于同一 Core，其间数据仍经 DDR，因而不能以``同核切分没有通信成本''作为模型假设。

设 Task \(s\) 在 Core \(k\) 上开始于 \(A_s\)、完成于 \(F_s\)，其同核前一个 Task 为 \(\mathrm{prev}(s)\)。激活时刻至少满足

\[A_s\ge F_{\mathrm{prev}(s)}+100,\]

\[A_s\ge F_t+1000,\quad t\in\operatorname{pred}(s),\ a(t)\ne a(s).\]

不存在的前项不参与最大值。这里 \(F_s-A_s\) 仍受其他核的 DDR 活动影响，不是只由子图结点周期相加得到的固定时长。细粒度切分增加可并行结构，粗粒度切分减少屏障和复制，两者之间的取舍构成问题一的核心。

算法首先保留依赖紧密、复用较强的局部结构，再为负载过重或前沿并行度不足的部分构造补充方案。这样的切分有明确对象：完整组件装箱解决独立工作分配；重组件和汇合尾部处理解决单一大块拖尾；分叉前沿与波内援助解决尚未暴露的并行分支。构造规则运行前冻结，由图结构触发，不从历史成绩表挑选方案。

\subsection{固定候选族及适用性守卫}\label{ux56faux5b9aux5019ux9009ux65cfux53caux9002ux7528ux6027ux5b88ux536b}

基础构造最多产生六种去重候选，包括组件有界装箱、重组件或汇合尾处理、过载分解、共享输入处理、分叉前沿和容量受限的私有返回链构造。共享输入规则考虑外部输入字节及组件关系；返回链规则先识别私有的 \(M\to V^+\to M\) 形状及容量条件。结构不满足时保留既有候选，不强行将一般 DAG 解释成链或树。

基础候选之上有三类有守卫的细化：容量返回赢家的变长分包响应细化；两种保持局部结构完整的前沿次序调整；波内三片援助。它们共享当前输入生成的父方案及其已知评分，以严格改善为接受条件。六个基础候选、至多一个响应细化、至多两个完整前沿候选和至多一个援助候选，使整个固定入口的在线 E1 评分上限为十次；去重、单候选快速路径和不适用分支使实际次数通常更少。

E1 用于在线完整方案比较；最终论文结果由官方 E0 的实际原件确认。响应编译保留核内 FIFO、依赖、存储门控、COPY 工作量、整数周期及共享带宽的关键结构。相同编译签名、共同空状态释放等条件下可以复用部分编译或响应计算，但这里不宣称一个有限签名已对任意未来输入充分，也不宣称有理数实现与官方 binary64 在全域严格等价。

\begin{figure}[!htbp]
\centering
\includegraphics[width=\linewidth,height=.84\textheight,keepaspectratio,trim=0 235 0 0,clip]{figures/fig-p1-cuts-zh.png}
\caption{case\_026、5 Core 的真实归约片段。左为原始划分，右为分叉援助；八个 ADD 节点的编号和依赖不变。点线端口连接裁剪范围外的输入或后继，橙色虚线为跨子图依赖。}
\label{fig:p1-cuts}
\end{figure}

\subsection{波内三片援助的构造与无环性}\label{ux6ce2ux5185ux4e09ux7247ux63f4ux52a9ux7684ux6784ux9020ux4e0eux65e0ux73afux6027}

在父方案的子图数据 DAG 中，令 \(h(s)\) 为 Task 的最长前驱高度。算法只处理每核顺序严格增加 \(h\) 的区域；同一高度、不同 Core 的 Task 之间必须不存在计算依赖。将负载较重核心称为 donor，较轻核心称为 helper，一次构造中二者角色互斥，helper 在同一波中须有可核对的原 Task，而不是任意空核。

在 donor 内选取符合识别条件的分叉---汇合结构，把计算划分为互不重叠的 \(X,Y,J\)：\(Y\) 为可导出的独立支路，\(X\) 为留在 donor 的其余前部，\(J\) 为汇合及后续尾部。要求 \(X,Y\) 之间没有依赖，\(J\) 不向 \(X,Y\) 返回，外部边与原波次方向相容。helper 先执行导入的 \(Y\)，再执行本波原 Task；donor 执行 \(X\) 后执行 \(J\)。构造按主导 Pipe 工作量选择负载转移，只有代理最大负载严格下降才继续形成完整候选。

\textbf{命题 4.1（三片援助的增广无环性）。} 在上述守卫条件下，并保持不同波按原高度先后排列，三片构造不在增广任务图中引入有向环。

\textbf{证明。} 给每个 Task 赋字典序秩 \((h,\phi,j)\)。第一项为原波高度；同波内，独立的 \(X\)、\(Y\) 放在前部阶段，helper 接收的多个 \(Y\) 按其实际顺序以 \(j\) 区分，helper 原 Task 位于其后，汇合尾 \(J\) 位于需要的 \(X,Y\) 之后。同波不同原 Task 没有计算边，\(X,Y\) 之间无边，\(J\) 无返回边，donor/helper 角色互斥，故所有新增数据边与同核序列边均沿同波秩向前。跨波边严格增加第一项。由命题 3.1，增广图无环。\(\square\)

证明只覆盖构造守卫对应的依赖与次序安全。导出支路会增加某些数据边界，新的 FIFO 与存储申请还可能改变换入换出。因此，代理负载下降不能作为接受收益的证据。构造完成后仍检查完整提交方案，并用完整 E1 结果比较 \((M,\text{scheduled COPY})\)。

\subsection{算法流程、停止条件与复杂度}\label{ux7b97ux6cd5ux6d41ux7a0bux505cux6b62ux6761ux4ef6ux4e0eux590dux6742ux5ea6}

\textbf{算法 1：屏障感知的固定结构候选求解。}

\begin{quote}
输入：原图 G，核数 K，冻结配置

1 解析图；按固定结构规则构造基础方案族，按规范化计划去重

2 若仅一份基础方案则直接保留；否则逐份完整评分，保持最佳已确认方案

3 若赢家满足容量返回识别条件，构造一次变长分包细化并比较

4 对冻结的两种完整前沿规则，逐一构造、检查、去重并比较

5 若父方案满足波次守卫，构造一次 X/Y/J 援助候选

6 检查完整覆盖、唯一调度及增广无环；候选严格改善时替换父方案

7 输出最佳已确认方案及独立诊断；诊断不写入提交 JSON
\end{quote}

解析、拓扑检查及给定候选的覆盖检查与图规模线性相关；排序、组件装箱和前沿优先队列增加相应的排序成本。三片构造还遍历可能的局部见证，不能将整个求解器简单标为一次线性遍历。该援助构造设置 Task 上限 320，援助结构检查的软时间预算为 2 s，附加评分最多一次；软预算通过检查点触发退出，并非每一步都可被精确抢占的硬时限。响应分包具有独立的受限编译预算和子进程停止条件，全部计入求解墙钟。

因完整入口最多十次 E1 评分，其评价成本为 \(\sum_{j=1}^{b}T_j\)、\(b\le10\)。图越大、存储事件越多，一次评分成本可能越高。实际实验同时保留每格调用数和进程收据，避免用``十次以内''替代实测效率。

\subsection{全量结果及质量---代价取舍}\label{ux5168ux91cfux7ed3ux679cux53caux8d28ux91cfux4ee3ux4ef7ux53d6ux820d}

固定算法在 100 图、1～5 核的 500 格中全部得到合法结果。题目要求的主曲线见下表；``优化单核''另列，避免改变官方主曲线的单核定义。

\begin{table}[!htbp]\centering\small
\begin{tabularx}{\linewidth}{@{}Xrrrrr@{}}\toprule
核数 & 1 & 2 & 3 & 4 & 5 \\
\midrule
官方平均加速比 & 1.0000 & 1.9516 & 2.7514 & 3.4786 & 4.0553 \\
优化方案 $\operatorname{mean}(B/M)$ & 1.0021 & 1.9516 & 2.7514 & 3.4786 & 4.0553 \\
改善/持平/退化 & 0/100/0 & 8/92/0 & 15/85/0 & 12/88/0 & 13/87/0 \\
求解墙钟均值 / s & 0.436 & 3.316 & 3.426 & 3.427 & 3.522 \\
\bottomrule\end{tabularx}\end{table}

5 核平均加速比由前版 4.0259 提升到 4.0553，相对增幅为 0.729\%。新三片援助在全部 500 格中被选中 36 次；其余格保留父方案。表中的完整版本比较包含固定策略组合，不能把所有改善都归因于三片援助的独立因果贡献。

5 核额外 DDR 总量为 1053684536 B，相比前版增加约 0.389\%。因此该次完整版本改进体现了以主要目标换取少量搬运增加的取舍，不能写成时间和 DDR 同时全面改善。5 核墙钟 P95 为 16.623 s，最大值为 25.552 s；该批次共享主机上使用四个 worker，操作系统缓存未清空，不能据此推断独占硬件的确定延迟。

全量调用账包含 500 次新求解进程、1130 次在线 E1，以及 48 次新 E0。另 452 格在原方案字节、图、核数、配置与评价语义逐格核同后复用已有 E0 原件；这些结果仍属于当前固定算法的输出，但旧 E0 墙钟不记作新测值。此区分使结果复用保持可追溯，也避免为了``全部新跑''重复消耗无新增信息的评价。

\begin{figure}[!htbp]
\centering
\includegraphics[width=\linewidth,height=.84\textheight,keepaspectratio]{figures/paper/fig-p1-overlap-old.pdf}
\caption{case\_026、5 Core 的父方案执行轨迹。上：Task 跨度；下：各 Core 的四条 Pipe。Makespan 为 44,114 cycles，额外搬运 211,680 B。}
\label{fig:p1-old}
\end{figure}

\begin{figure}[!htbp]
\centering
\includegraphics[width=\linewidth,height=.84\textheight,keepaspectratio]{figures/paper/fig-p1-overlap-new.pdf}
\caption{同一用例的分叉援助方案。横轴与父方案一致；Task 数从 5 增为 7，Makespan 降至 42,014 cycles，额外搬运增至 375,532 B。条带均来自官方执行事件。}
\label{fig:p1-new}
\end{figure}

\begin{figure}[!htbp]
\centering
\includegraphics[width=\linewidth,height=.84\textheight,keepaspectratio]{figures/paper/fig-p1-results.pdf}
\caption{问题一的完整版本比较。上：每个核数的 100 图平均加速比；下：全部 500 格的时间与搬运量配对变化。核数用点形区分，横轴采用对称对数。}
\label{fig:p1-results}
\end{figure}

本章说明，结构化切分可以在保留真实评价闭环的同时，减少盲目枚举。剩余瓶颈可能是不可拆分的计算链、边界同步、共享输入或动态存储行为；因此下一问利用同核驻留后，优化重点需要从``少设 Task 屏障''进一步转向``复用与发射次序共同安排''。

\FloatBarrier

\section{问题二：复用感知的超图分核与空隙插入}\label{ux95eeux9898ux4e8cux590dux7528ux611fux77e5ux7684ux8d85ux56feux5206ux6838ux4e0eux7a7aux9699ux63d2ux5165}

\subsection{同核合并改变了什么}\label{ux540cux6838ux5408ux5e76ux6539ux53d8ux4e86ux4ec0ux4e48}

问题二把同一 Core 上的子图合并为一个 Task。同核依赖可以保留核内数据，跨核依赖才插入边界 COPY，并在源 COPY\_OUT 完成后经过 500 cycles 才允许目标 COPY\_IN 进入相应就绪判定。于是，拆成 singleton 子图不再等于为每个操作设置独立 Task 屏障；子图顺序成为影响固定核内调度的重要输入。

这一变化同时带来机会和压力。把共享数据的操作放在同核可减少重复搬运，但数据驻留时间可能延长；按计算负载分散消费者可改善并行度，却扩大触及 Core 集合和跨核同步。本文分别构造计算日历代理与精确的预换入换出 COPY 费用，再由完整方案评价决定二者的最终取舍。

\subsection{链与汇合感知的空隙插入}\label{ux94feux4e0eux6c47ux5408ux611fux77e5ux7684ux7a7aux9699ux63d2ux5165}

对计算依赖图识别可连续处理的链以及分叉、汇合位置。每个 Core 的 M、V Pipe 维护已安排区间集合，候选操作必须在依赖释放后寻找足够长的空隙。相比仅在某核当前尾部追加工作，空隙插入可以利用前序等待留下的可用窗口。

设操作 \(v\) 的代理释放时刻为 \(r_v\)，计算时长为 \(d_v\)。对 Core \(k\)、Pipe \(p_v\) 的占用日历 \(\mathcal I_{k,p_v}\)，选择最早满足

\[[s_v,s_v+d_v)\cap\mathcal I_{k,p_v}=\varnothing,\qquad s_v\ge r_v\]

的区间。汇合附近的相关链使用至多 \(K^2\) 个核对联合比较，使一个支路的局部选择考虑到另一个支路的完成位置。输出以 singleton 子图及其优先序表达，后续仍由官方规则生成实际流水。

这个日历模型故意保持可计算：它不完整模拟共享 DDR、公平分带宽、spill 与 memory-credit。\(s_v\) 是构造代理时刻，不是官方 trace，不是 Makespan 的保守上下界。日历的价值在于产生结构合理的候选，而不是取代官方执行。

同核合并与跨核复制的区别见图\textasciitilde{}\ref{fig:progression}。后续分核模型据此区分复用收益与跨核费用。

\subsection{按触及 Core 数计费的超图模型}\label{ux6309ux89e6ux53ca-core-ux6570ux8ba1ux8d39ux7684ux8d85ux56feux6a21ux578b}

普通图切割按边累加费用，容易对同一张量的多消费者重复计费。超图用一条超边连接同一数据相关的全部操作，以触及 Core 数描述通信，这种建模思想与并行稀疏矩阵计算中的通信建模相通{[}3{]}。本文根据官方边界复制语义构造具体费用。

要求非 COPY 操作集合完整、物理张量不带逻辑别名、每个物理张量至多有一个原生产者，张量大小为非负整数。对超边 \(e\)，令 \(\lambda_e(a)\) 为其端点所触及的不同 Core 数，\(w_e\) 为字节权重。则预 Step2 的原始 COPY 费用写成

\[D_{\mathrm{pre2}}(a)=C+\sum_{e\in\mathcal H}w_e\bigl(\lambda_e(a)-1\bigr).\]

对于核外输入，消费者每增加一个触及的 Core，就增加一次大小为 \(b_t\) 的输入复制，故权重为 \(b_t\)；基础的一次输入包含在常数 \(C\) 中。对于由计算产生的中间张量，超边包含生产者和消费者，官方边界构造按远端核所需的输出/输入对计费，权重为 \(2b_t\)；最终输出等与分核无关的部分计入 \(C\)。实际实现按官方边界规则处理，不能把所有边一律赋为相同权重。

该公式精确描述适用域内的预 Step2 COPY 字节，不包含容量不足引发的换入换出，也不表达传输之间的重叠或串行。若一组操作 \(A\) 从 Core \(u\) 移往 \(v\)，只需检查接触 \(A\) 的超边。对某条超边，目标核原先无端点时费用增加一个 \(w_e\)；源核全部端点都迁出时费用减少一个 \(w_e\)。由此可增量计算

\[\Delta_e=w_e\left(\mathbf1\{n_{e,v}=0\}-\mathbf1\{n_{e,u}=|e\cap A|\}\right).\]

这也解释了单结点逐次迁移可能看不到组合收益：第一步迁移增加目标核触及数，只有最后一个源端点迁出时才消除源核费用。局部联合改动比纯逐点贪心更适合保留复用结构。

以一个张量的超边为例，令 \(u\) 为生产者，\(v_1,v_2\) 为两个消费者，并以 \(t\) 标记该张量对应的超边。将两个消费者中的一个迁入生产者所在 Core 后，另一个消费者仍使超边触及两个 Core，费用项暂不下降；只有两个消费者都迁入时，\(\lambda_t\) 才从 2 降为 1，对应费用项从 \(w_t\) 降为 0。图中的变化仅针对该超边，不能解释为完整方案的 COPY 字节全部归零。

\begin{figure}[!htbp]
\centering
\includegraphics[width=\linewidth,height=.84\textheight,keepaspectratio]{figures/fig-hypercut-paper.png}
\caption{张量超边的触核数与消费者联合迁移示意。左：生产者与消费者分处两个 Core；右：两个消费者均迁入生产者所在 Core。橙色连线表示超边关联。}
\label{fig:hypercut}
\end{figure}

这一局部字节收益仍需接受执行层的检验：联合迁移可能同时改变计算负载、数据驻留和同步位置，而上述费用没有刻画共享带宽、容量换入换出及流水重叠。下一步因此对局部割加入负载守卫，再用完整方案评价判断 Makespan 是否改善。

\subsection{带锚点的两核局部细化与负载守卫}\label{ux5e26ux951aux70b9ux7684ux4e24ux6838ux5c40ux90e8ux7ec6ux5316ux4e0eux8d1fux8f7dux5b88ux536b}

在已构造的 gap 方案上，算法选择核对与最多 16 条链组成的局部区域。区域之外的操作及必须固定的链作为锚点，区域内操作具有两种 Core 标签。对固定区域和锚点，超边``是否同时触及两侧''的费用可通过辅助结点构造成二标签最小割问题，得到该局部字节目标的精确解。

但无约束最小割可能将多数工作迁入同一 Core，造成计算负载失衡。实现为各核各 Pipe 设置工作帽；若新分配违反工作帽，逐步将部分迁入链锚回原核后重新求割。原分配始终是可行参考，最多 \(r+1\) 次割（\(r\le16\)）后得到满足守卫的接受候选或保留原分配。每次接受还要求字节数不增加。

\textbf{命题 5.1（局部费用单调性）。} 若每轮锚点均保持原分配可行，局部割求得其当前无额外工作帽子问题的精确最优值，且接受前重新核对工作帽，则最终接受方案在该局部预 Step2 COPY 目标上不劣于原分配。

\textbf{证明。} 当前锚点约束下原分配可行，因此精确最小割值不大于原分配值。若工作帽不满足，则增加与原分配相容的锚点，不改变原分配可行性。每轮仍保留这一费用上界。只有同时通过工作帽和实际费用检查的候选才接受，结论成立。\(\square\)

该过程不等于带所有负载约束的全局精确最小割。新增锚点会缩小可行域，也未纳入容量、共享带宽或真实 Makespan。本文将其称为带守卫的局部字节细化，避免把局部精确性误写为多核调度全局最优。

\subsection{完整候选策略与复杂度}\label{ux5b8cux6574ux5019ux9009ux7b56ux7565ux4e0eux590dux6742ux5ea6}

\textbf{算法 2：复用感知的有界完整方案选择。}

\begin{quote}
输入：G、K、冻结配置及完整方案评分接口

1 生成结构自适应基线 P0

2 构造链/汇合感知 gap 方案 P1；不适用时返回 P0

3 在 P1 上做两核、至多16链的超图细化

4 仅当预 Step2 COPY 严格下降时，对细化分配重新安排次序，形成 P2

5 保留未经细化的 P1；按 P0、P1、P2 的构造顺序去重

6 对至多三个完整计划在线评分，按 (Makespan, added COPY) 字典序择优

7 同分保留先前方案；失败按固定回退规则处理；输出合法两字段 JSON
\end{quote}

保留 \(P_1\) 很关键：预 Step2 字节没有下降，只能说明不触发细化后重排，不能据此删除原 gap 方案。在线最终选择并不以字节代理代替 Makespan。

设图中有 \(n\) 个计算操作、\(m\) 条投影依赖，日历查询累计扫描成本为 \(Q\)；链、汇合识别与相应排序的成本由 \(n,m\) 决定。若有 \(R\) 个局部细化区域，区域辅助网络规模为 \((n_r,m_r)\)，一次最小割成本为 \(T_{\mathrm{cut}}(n_r,m_r)\)，则细化总成本不超过

\[\sum_{r=1}^{R}(w_r+1)T_{\mathrm{cut}}(n_r,m_r),\qquad w_r\le16.\]

日历插入及核对比较还包含 \(Q\) 与 \(K^2\) 因子，因此不将其笼统称为线性算法。最昂贵的完整方案在线评分最多三次，区域局部割不调用完整评价器。各阶段时间均计入求解墙钟。

\subsection{完整结果与通信---时间取舍}\label{ux5b8cux6574ux7ed3ux679cux4e0eux901aux4fe1ux65f6ux95f4ux53d6ux820d}

冻结求解器在全部 500 格上输出合法结果，使用 1131 次原生 E2 评分和 500 次独立 E0 复评，未触发 E0 fallback。主结果为：

\begin{table}[!htbp]\centering\small
\begin{tabularx}{\linewidth}{@{}Xrrrrr@{}}\toprule
核数 & 1 & 2 & 3 & 4 & 5 \\
\midrule
官方平均加速比 & 1.0000 & 2.3167 & 3.2356 & 3.9711 & 4.5498 \\
优化方案 $\operatorname{mean}(B/M)$ & 1.2061 & 2.3167 & 3.2356 & 3.9711 & 4.5498 \\
改善/持平/退化 & 59/41/0 & 64/36/0 & 54/46/0 & 49/51/0 & 42/58/0 \\
\bottomrule\end{tabularx}\end{table}

相对前一完整版本，500 格共有 268 格 Makespan 改善、232 格持平、0 格退化，涉及至少一个核数改善的图有 69 个。这个结果证明固定组合策略在本组已知官方图上的经验表现，不能证明每个组成机制对所有输入独立保优。

全部 500 格的额外 DDR 总量由 3043860416 B 增至 4159738560 B，增加约 36.66\%。其中 55 格时间与搬运均改善，199 格时间改善而搬运增加，14 格时间改善且搬运不变，232 格两者持平。更大的总搬运伴随更短 Makespan，表明仅最小化通信字节不足以描述本题主要目标。

求解墙钟均值为 4.029 s，最近秩定义的 P95 为 18.841 s，最大值为 40.187 s；外部独立 E0 的均值为 1.410 s、P95 为 6.592 s、最大值为 13.980 s。实验使用共享 Mac、单 worker，整批约 46.63 min。此处的``新进程''不意味着清空操作系统磁盘缓存或获得独占算力。

\begin{figure}[!htbp]
\centering
\includegraphics[width=\linewidth,height=.84\textheight,keepaspectratio]{figures/paper/fig-p2-results.pdf}
\caption{问题二的完整版本比较。上：官方口径的平均加速比；下：500 格配对结果。部分时间改善伴随额外搬运增加，两项指标分别报告。}
\label{fig:p2-results}
\end{figure}

\subsection{负例：零 spill 与少搬运仍不足以保证更快}\label{ux8d1fux4f8bux96f6-spill-ux4e0eux5c11ux642cux8fd0ux4ecdux4e0dux8db3ux4ee5ux4fddux8bc1ux66f4ux5feb}

为检验通信与容量代理的解释力，另一个冻结的零 spill 联合构造 C04 在六个预选 5 核用例上接受官方评价。六格都合法，零 spill 及字节核对成立，但 Makespan 全部劣于本章主算法。

{\def\LTcaptype{table} % do not increment counter
\begin{longtable}[]{@{}lrrr@{}}
\toprule\noalign{}
case & 主算法 / cycles & C04 / cycles & 相对变慢 \\
\midrule\noalign{}
\endhead
\bottomrule\noalign{}
\endlastfoot
005 & 33515 & 38209 & 14.01\% \\
016 & 2091375 & 2399125 & 14.72\% \\
019 & 16247 & 28514 & 75.50\% \\
025 & 942938 & 975658 & 3.47\% \\
039 & 449653 & 494942 & 10.07\% \\
075 & 343930 & 452120 & 31.46\% \\
\end{longtable}
}

其中 019、039、075 的 COPY 字节下降而完成时间上升。原因不能仅凭这两个聚合量确定；该结果已经足以否定``预估字节降低且零 spill 就必然更快''的推论。计算 FIFO 不变也未必保留 COPY FIFO、共享 DDR 竞争和跨核释放过程。由此，零 spill 证书若只覆盖官方 Step2，不应被宣传为完整执行最优证书。

\begin{figure}[!htbp]
\centering
\includegraphics[width=\linewidth,height=.84\textheight,keepaspectratio]{figures/paper/fig-p2-counterexample-control.pdf}
\caption{case\_019、5 Core 主算法的真实执行时间线。上为合并后的 Task，下为四条 Pipe；Makespan 为 16,247 cycles。}
\label{fig:p2-control}
\end{figure}

\begin{figure}[!htbp]
\centering
\includegraphics[width=\linewidth,height=.84\textheight,keepaspectratio]{figures/paper/fig-p2-counterexample-c04.pdf}
\caption{同一用例的 C04 候选，使用相同横轴。搬运量下降且 spill 为零，但 Makespan 增至 28,514 cycles，说明字节代理不足以判定执行收益。}
\label{fig:p2-c04}
\end{figure}

与问题一相比，本章把同核驻留带来的收益纳入候选构造，同时保留对动态执行的谨慎。下一章再加入共享 Cache 时，复用机会还要从``消费者在哪个 Core''延伸到``消费者何时请求''。

\FloatBarrier

\section{问题三：归约森林次序与只读 Cache 协同}\label{ux95eeux9898ux4e09ux5f52ux7ea6ux68eeux6797ux6b21ux5e8fux4e0eux53eaux8bfb-cache-ux534fux540c}

\subsection{从空间复用到请求次序}\label{ux4eceux7a7aux95f4ux590dux7528ux5230ux8bf7ux6c42ux6b21ux5e8f}

问题三保留问题二的同核合并和跨核 COPY 语义，增加所有核心共享的只读 Cache。COPY\_IN 发射时按逻辑张量标识查询：命中进入 Cache 读带宽池；未命中进入 DDR 池，传输完成后填入 Cache。命中不改变 FIFO 次序；容量不足时按进入顺序淘汰。

Cache 不消除原数据依赖，不免除源 COPY\_OUT 完成后的 500 cycles 跨核等待，也不把命中读变成零时长操作。它改变的是读请求使用的服务资源与后续状态。故 \(M_3(P)\) 取决于请求序列、发射时刻、共享输入和在途传输，不能仅用总命中字节乘一个固定节省系数估算。

本文沿用前两问的结构---评价组织方式：首先按图形态构造合法的分核和序列；对共享输入流水、注意力结构中的次序及私有归约森林分别提供有守卫候选；再以固定候选下界排除不可能改善者，最终只接受完整官方评价的严格 Makespan 改善。主算法称为 Forest，指包含森林次序候选的完整固定策略，而不意味着所有输入都是归约树。

例如，case\_021、3 Core 的编译后张量 t1000000001 大小为 4096 B。官方事件给出以下生命周期；事件时刻为 cycles，表中间隔不表示等长处理阶段。

\begin{table}[!htbp]\centering
\caption{同一张量的 Cache 事件}\label{tab:cache-events}
\begin{tabularx}{\linewidth}{rX}\toprule
时刻 / cycle & 事件 \\\midrule
0 & COPY\_IN 发射时未命中，尚未填充 \\
207 & DDR COPY\_IN 完成，张量插入 FIFO \\
4,956 & 另一 Core 读命中，FIFO 年龄不刷新 \\
323,850 & 插入张量 t1000000487 时，该最老项被淘汰 \\\bottomrule
\end{tabularx}\end{table}

\subsection{私有归约树的内部存活前沿}\label{ux79c1ux6709ux5f52ux7ea6ux6811ux7684ux5185ux90e8ux5b58ux6d3bux524dux6cbf}

归约树的多个子树先产生结果，父操作等子结果齐备后消费。处理子树的先后顺序决定先完成的结果需要保留多久。在固定私有树、每个内部结果唯一生产且只由父结点消费、每结点产生一个已知大小输出的条件下，可以把子树压缩为一对量：内部峰值 \(h_i\) 与最终保留结果大小 \(r_i\)。

此处限定非交错执行：一个子树完成后才处理下一个兄弟；只统计内部结果的标量字节；父输出先分配，再释放孩子结果。令兄弟次序为 \(\pi\)，则峰值为

\[H(\pi)=\max\left\{\max_j\left(\sum_{\ell<j}r_{\pi_\ell}+h_{\pi_j}\right),\ \sum_i r_i+r_{\mathrm{parent}}\right\}.\]

第一项表示处理第 \(j\) 个子树时，之前兄弟留下的结果与当前子树内部峰值叠加；第二项表示父输出分配瞬间仍须保留所有孩子结果。这个式子给出了``先做哪个分支''的具体数学对象。

\textbf{命题 6.1（内部前沿排序）。} 若 \(h_i\ge r_i\ge0\)，按 \(h_i-r_i\) 非增顺序处理兄弟子树，在上述非交错模型中最小化 \(H\)。

\textbf{证明。} 考虑相邻的 \(A,B\)，且 \(h_A-r_A\ge h_B-r_B\)。顺序 \(A,B\) 与 \(B,A\) 的局部峰值分别为

\[H_{AB}=\max(h_A,r_A+h_B),\quad H_{BA}=\max(h_B,r_B+h_A).\]

由条件 \(r_A+h_B\le r_B+h_A\)，且 \(h_A\le r_B+h_A\)，故 \(H_{AB}\le H_{BA}\)。两种次序处理后均保留 \(r_A+r_B\)，不会改变后续兄弟的初始存活量；共同前缀只为二者峰值加上同一个常量，父分配项也不变。因此交换逆序相邻对不增加峰值。对兄弟次序重复交换，再对子树自底向上归纳，得到结论。\(\square\)

这一交换论证与经典表达式树代码生成中的存储次序思想相关{[}4{]}，本文的适配对象是已知字节大小的内部结果。算法据此递推 \(h,r\)，生成深度优先后序，保持原来的归核结果，再用 singleton 子图序列表达。森林识别还要求原始 COPY 只在合法边界出现、内部输出无 fanout、无缺少张量大小的额外直接计算边等。

命题最小化的是内部前沿代理。外部输入、不同归约树之间的根结果保留、L1/UB 双池、编译后物理实例、COPY 和真实存储申请次序不在这个标量模型内。因此它不是全局缓存峰值最优，也不是零 spill 或 Makespan 最优定理。官方评价是将有理论依据的次序转化为机器层结论的必要一步。

取 \(h_A=8,r_A=2,h_B=7,r_B=5\)，父输出大小为 \(3\)（单位均为 KiB）。按 \(A\to B\) 处理时，局部峰值为 \(\max\{8,2+7,2+5+3\}=10\)；反序峰值为 \(\max\{7,5+8,5+2+3\}=13\)。该例保留父输出先分配的费用，只解释内部前沿排序，不是官方存储池的实测轨迹。

\subsection{固定候选的下界剪枝与接受规则}\label{ux56faux5b9aux5019ux9009ux7684ux4e0bux754cux526aux679dux4e0eux63a5ux53d7ux89c4ux5219}

在 singleton 计划、Core 归属和同核同 Pipe FIFO 都已经确定时，将有证据的必要计算依赖、FIFO 先后及跨核释放约束构成加权有向图。对每条必要关系 \(u\to v\) 写成

\[s_v\ge s_u+d_u+\delta_{uv},\]

其中 \(\delta_{uv}\) 只包含已经证明不可避免的等待；同端点多条约束取最大值，不相加。图上最长路径 \(L(P)\) 是固定候选任意成功执行的下界。若当前 fresh incumbent 的官方完成时间为 \(U\)，则 \(L(P)\ge U\) 时，该候选不可能严格改善，可省去完整 E0。

这里不能使用任意``去掉资源限制后重新贪心调度''的耗时作为下界：资源放松扩大了可行集合，却不保证一个固定贪心策略的输出单调改善。本文的下界来自每条真实执行都满足的不等式与路径，而不是放松模拟器恰好给出的一个解。

\textbf{算法 3：Forest 完整策略。}

\begin{quote}
输入：G、K、冻结配置

1 从当前输入建立结构索引，生成并评价结构/日历基准方案

2 对适用结构，考虑共享输入连续分段或保留注意力分配见证的次序

3 在剩余总预算内，识别私有归约森林并生成 h-r 排序候选

4 对重复计划直接跳过；对满足条件的候选计算固定计划下界

5 若下界不可能击败 incumbent，则保留剪枝证据并跳过 E0

6 否则作完整官方 P3 评价；只有 Makespan 严格降低才替换

7 总在线 E0 不超过3次，输出最佳成功方案及完整调用记录
\end{quote}

三个候选名额为整次调用共享，不是每个机制各有三次。某些结构的早期候选已使用预算时，后续森林候选不会继续新增评价。结构不适用、重复、下界剪枝、合法性拒绝及实际收益不佳是不同停止原因。

对给定归约森林，识别和后序遍历为 \(O(n+m)\)；兄弟排序为 \(O(\sum_v\deg(v)\log\deg(v))\)。固定候选 DAG 的最长路径为 \(O(n+m')\)，其中 \(m'\) 包含必要 FIFO 等边；图不在已证域内时不使用该下界。整体还包含结构索引、分配构造与至多三次官方评价，不能将整个 Forest 的复杂度等同于树排序子过程。

\subsection{同计划 Cache 对照与完整算法结果}\label{ux540cux8ba1ux5212-cache-ux5bf9ux7167ux4e0eux5b8cux6574ux7b97ux6cd5ux7ed3ux679c}

Forest 固定版本在 100 图×1～5 核上全部得到官方成功结果，500 次新求解过程共使用 973 次在线 E0，无重试。相对前一完整策略，24 格 Makespan 改善，476 格持平，0 格退化。额外 COPY 与 spill 各有 18 格减少、479 格持平、3 格增加，净变化各为 −413450240 B；spill 是总搬运的组成部分，二者不能相加声称减少了双倍字节。

为了隔离 Cache 作用，对每个最终计划使用相同图、核数与计划分别评价问题二和问题三。这里的无 L2 对照是 \textbf{P3 计划在 P2 语义下的结果}，不是问题二主算法另行优化出的赢家。

{\def\LTcaptype{table} % do not increment counter
\begin{longtable}[]{@{}
  >{\raggedright\arraybackslash}p{(\linewidth - 8\tabcolsep) * \real{0.1579}}
  >{\raggedleft\arraybackslash}p{(\linewidth - 8\tabcolsep) * \real{0.2105}}
  >{\raggedleft\arraybackslash}p{(\linewidth - 8\tabcolsep) * \real{0.2105}}
  >{\raggedleft\arraybackslash}p{(\linewidth - 8\tabcolsep) * \real{0.2105}}
  >{\raggedleft\arraybackslash}p{(\linewidth - 8\tabcolsep) * \real{0.2105}}@{}}
\toprule\noalign{}
\begin{minipage}[b]{\linewidth}\raggedright
核数
\end{minipage} & \begin{minipage}[b]{\linewidth}\raggedleft
无 L2：\(\operatorname{mean}(B/M_2)\)
\end{minipage} & \begin{minipage}[b]{\linewidth}\raggedleft
Cache：\(\operatorname{mean}(B/M_3)\)
\end{minipage} & \begin{minipage}[b]{\linewidth}\raggedleft
\(\operatorname{mean}(M_2/M_3)\)
\end{minipage} & \begin{minipage}[b]{\linewidth}\raggedleft
逐图字节命中率均值
\end{minipage} \\
\midrule\noalign{}
\endhead
\bottomrule\noalign{}
\endlastfoot
1 & 1.1971 & 1.1998 & 1.0023 & 6.47\% \\
2 & 2.2949 & 2.3120 & 1.0079 & 16.25\% \\
3 & 3.1745 & 3.2587 & 1.0314 & 25.96\% \\
4 & 3.9211 & 4.0915 & 1.0512 & 24.53\% \\
5 & 4.4763 & 4.7576 & 1.0829 & 31.36\% \\
\end{longtable}
}

5 核的逐图 Cache 加速比平均为 1.0829，但这不是``所有用例都快了 8.29\%''，也不是完成时间百分比下降的平均值。跨图访问字节加权的 5 核命中率为 47.45\%，与上表的逐图算术平均 31.36\% 含义不同，大流量图在前者中权重更大。

Forest 求解墙钟均值为 4.401 s、中位数 1.616 s、最近秩 P95 为 20.009 s、最大值为 51.335 s。它包括在线 E0，来源是共享主机四 worker 的完整批次。后续 R9F 批次的单 worker 时间不用于替换这些数值。

\begin{figure}[!htbp]
\centering
\includegraphics[width=\linewidth,height=.84\textheight,keepaspectratio]{figures/paper/fig-p3-results.pdf}
\caption{问题三的同计划配对评价。上：无 L2 与只读 Cache 的平均基线加速比；下：逐图 $M_2/M_3$ 的均值。每个核数包含 100 对，不用均值之比替代比值均值。}
\label{fig:p3-results}
\end{figure}

\begin{figure}[!htbp]
\centering
\includegraphics[width=\linewidth,height=.84\textheight,keepaspectratio]{figures/paper/fig-cache-gain.pdf}
\caption{500 个同计划配对结果中的字节命中率与 Cache 收益。保留低于 1 的负收益点；100 个用例的不同核数结果相关，不能当作 500 个独立随机样本。}
\label{fig:cache-gain}
\end{figure}

\subsection{选择政策的负消融：相对收益与绝对收益可能冲突}\label{ux9009ux62e9ux653fux7b56ux7684ux8d1fux6d88ux878dux76f8ux5bf9ux6536ux76caux4e0eux7eddux5bf9ux6536ux76caux53efux80fdux51b2ux7a81}

后续 R9F 对 5 核构造增加了一项联合门槛：候选须使 P3 的 \(M_3\) 严格降低、同计划 P2 的 \(M_2\) 不增加，同时 \(G=M_2/M_3\) 不降低。100 个 5 核图中，96 个不满足构造识别条件，1 个由候选下界剪枝，3 个被联合门槛拒绝，最终接受数为零。全 500 个最终方案均与 Forest 字节相同。

三个被拒绝的候选实际上同时降低了无 L2 和 Cache 的绝对完成时间：

{\def\LTcaptype{table} % do not increment counter
\begin{longtable}[]{@{}lrrrr@{}}
\toprule\noalign{}
case/K & 原 \(M_2\) & 候选 \(M_2\) & 原 \(M_3\) & 候选 \(M_3\) \\
\midrule\noalign{}
\endhead
\bottomrule\noalign{}
\endlastfoot
064/5 & 7029 & 6899 & 6981 & 6877 \\
068/5 & 131631 & 98913 & 116345 & 97971 \\
088/5 & 97305 & 66729 & 85789 & 66224 \\
\end{longtable}
}

联合门槛中的比值条件等价于

\[M_{2,\mathrm{new}}M_{3,\mathrm{old}}\ge M_{2,\mathrm{old}}M_{3,\mathrm{new}}.\]

由于分子随方案改变，无 L2 时间改善得更快时，Cache 相对收益可能下降，即使 \(M_3\) 已明显缩短。因此这个负消融表明``维持相对 Cache 收益''的政策可能排除主要目标更优的方案，不能解释成第五核构造完全没有绝对收益。该门槛是研究时采用的选择政策，官方首要目标仍是 Makespan。

本文仍以 Forest 的固定全量作为主结果，并明确它不是所有已测单格候选中的最小 \(M_3\)。将三格候选手工替换到主表会变成事后逐格选择，不能代表一个已经完成全量验证的新算法。R9F 的 400+100 两段独立执行共计982次在线E0，只用于说明这项冻结政策的实际行为。

本章把计算树上的可证明排序、候选级必要条件和共享 Cache 的真实观测放在同一条论证链上。最有价值的结论既包括何时能减少内部前沿，也包括这种局部最优性为什么不足以取代最终执行评价。

\FloatBarrier

\section{实验设计与综合分析}\label{ux5b9eux9a8cux8bbeux8ba1ux4e0eux7efcux5408ux5206ux6790}

\subsection{数据、版本与对照协议}\label{ux6570ux636eux7248ux672cux4e0eux5bf9ux7167ux534fux8bae}

实验采用赛题给定的全部100个原始计算图，核数取1～5。每题对应500格，三题共1500格。所有主成绩使用同一题内冻结的完整求解器、统一入口、固定候选规则和预算，不以每个格子的历史最优组合代表一个算法。原始图和配置不清洗、不改写；逐图基线与对应的官方结果绑定。

本文区分三种比较。第一种是官方单核启发式基线与多核方案之比，用于题目要求的加速曲线；第二种是同题、同图、同核数下两个完整固定版本的结果，用于观察方法演进；第三种是相同P3方案在无L2与只读Cache条件下的结果，用于隔离Cache作用。三问硬件语义本身不同，不用P1与P2曲线的差直接估计某一个算法模块的贡献。

正文表格和后续图件使用同一份逐格结果表。导出程序对100×5覆盖、单核分母身份及原件字节进行核对，对P3另外读取同计划无L2原件，重算每核平均值。独立作者审计提供完整计划、官方结果及调用账的对应证据；论文阶段的数据导出是静态复核，不声称重新执行了全部官方评分。

\begin{figure}[!htbp]
\centering
\includegraphics[width=\linewidth,height=.84\textheight,keepaspectratio]{figures/paper/fig-dataset.pdf}
\caption{官方 100 张计算图的结构统计。上：非 COPY 操作数与张量大小总和；下：按矩阵 Pipe 工作占比排序的用例分布。}
\label{fig:dataset}
\end{figure}

\subsection{主要质量结果}\label{ux4e3bux8981ux8d28ux91cfux7ed3ux679c}

下表为各问题优化方案相对共同官方单核分母的逐图比值均值。P1/P2题目主曲线的1核坐标另按规定设为1；本表保留算法优化单核表现，以揭示基线本身的启发式性质。

{\def\LTcaptype{table} % do not increment counter
\begin{longtable}[]{@{}lrrrr@{}}
\toprule\noalign{}
核数 & P1优化方案 & P2优化方案 & P3优化方案 & P3同计划无L2 \\
\midrule\noalign{}
\endhead
\bottomrule\noalign{}
\endlastfoot
1 & 1.0021 & 1.2061 & 1.1998 & 1.1971 \\
2 & 1.9516 & 2.3167 & 2.3120 & 2.2949 \\
3 & 2.7514 & 3.2356 & 3.2587 & 3.1745 \\
4 & 3.4786 & 3.9711 & 4.0915 & 3.9211 \\
5 & 4.0553 & 4.5498 & 4.7576 & 4.4763 \\
\end{longtable}
}

P1的完整Task边界限制了跨子图复用；P2/P3允许同核合并，但它们仍受共享DDR、跨核释放和缓存压力影响。表中的结果与这个递进关系相符，却不能单独证明其差异全部由某一项资源解释。机制证据需要回到同条件方案对照及事件轨迹。

此外，P2主算法与``P3同计划无L2''一列是两个不同的方案生成策略。后者只是为P3隔离Cache作用的对照，不应被解释为P2主算法退化。5核P3的4.7576与无L2的4.4763之比约为1.0628，而逐图\(M_2/M_3\)的均值为1.0829；不同聚合顺序产生不同结果，正文以题目要求和配对定义为准。

\subsection{求解效率与资源成本}\label{ux6c42ux89e3ux6548ux7387ux4e0eux8d44ux6e90ux6210ux672c}

三种方法均把结构预处理、候选构造、在线评分及必要输出包含在求解墙钟中。P1/P2在线使用加速评价接口，最终E0复评独立计时；P3将官方E0直接用于在线选择，其时间已经包含在求解器内。

{\def\LTcaptype{table} % do not increment counter
\begin{longtable}[]{@{}
  >{\raggedright\arraybackslash}p{(\linewidth - 10\tabcolsep) * \real{0.1364}}
  >{\raggedleft\arraybackslash}p{(\linewidth - 10\tabcolsep) * \real{0.1818}}
  >{\raggedleft\arraybackslash}p{(\linewidth - 10\tabcolsep) * \real{0.1818}}
  >{\raggedleft\arraybackslash}p{(\linewidth - 10\tabcolsep) * \real{0.1818}}
  >{\raggedleft\arraybackslash}p{(\linewidth - 10\tabcolsep) * \real{0.1818}}
  >{\raggedright\arraybackslash}p{(\linewidth - 10\tabcolsep) * \real{0.1364}}@{}}
\toprule\noalign{}
\begin{minipage}[b]{\linewidth}\raggedright
方法与样本范围
\end{minipage} & \begin{minipage}[b]{\linewidth}\raggedleft
在线评价总次数
\end{minipage} & \begin{minipage}[b]{\linewidth}\raggedleft
求解均值 / s
\end{minipage} & \begin{minipage}[b]{\linewidth}\raggedleft
P95 / s
\end{minipage} & \begin{minipage}[b]{\linewidth}\raggedleft
最大 / s
\end{minipage} & \begin{minipage}[b]{\linewidth}\raggedright
运行条件
\end{minipage} \\
\midrule\noalign{}
\endhead
\bottomrule\noalign{}
\endlastfoot
P1，100个5核格 & 全500格1130次E1 & 3.522 & 16.623 & 25.552 & 共享主机，4 worker \\
P2，全部500格 & 1131次E2 & 4.029 & 18.841 & 40.187 & 共享主机，1 worker \\
P3 Forest，全部500格 & 973次E0 & 4.401 & 20.009 & 51.335 & 共享主机，4 worker \\
\end{longtable}
}

P1作者表的P95采用线性插值定义，P2/P3采用最近秩定义；不同定义在图表中注明，后续统一分布图可从逐格时间重新计算同一口径，不能混用后宣称精细排名。上表P1的时间只覆盖5核，评价次数覆盖全500格，两个范围特意分开标注。

题面建议单个用例在5～10分钟内给出方案，上述成功进程的实测时间明显小于这一量级。但实验未进行独占硬件、清空操作系统缓存或跨平台重复测量，故只能报告给定共享环境的测量范围。算法设计仍以同时降低官方M和端到端墙钟为目标，不以进入10分钟以内作为效率优化的终点。

\begin{figure}[!htbp]
\centering
\includegraphics[width=\linewidth,height=.84\textheight,keepaspectratio]{figures/paper/fig-runtime.pdf}
\caption{三种完整求解器的端到端墙钟经验分布，每题 500 次。时间含各自在线评价，横轴为对数；共享主机和并发条件不同，不能据此作独占硬件速度排序。}
\label{fig:runtime}
\end{figure}

\subsection{机制证据、消融与负结果}\label{ux673aux5236ux8bc1ux636eux6d88ux878dux4e0eux8d1fux7ed3ux679c}

三类已有证据承担不同任务。P1/P2/P3前后完整版本对照展示最终固定策略的整体变化；每个被选中机制的记录说明它确实触发并影响了输出；真实局部结构和时间线解释一格中发生了什么。若多个机制同时变化，前后版本差不能自动分摊成各模块的独立贡献。

P1三片援助提供了新增并行机会，也增加了部分边界搬运；5核最终版本的平均加速比改善约0.729\%，额外DDR总量增加约0.389\%。P2完整新版本在268格上缩短M，但总额外DDR增加约36.66\%；这一取舍说明并行与流水收益可以超过新增通信的代价。P3森林次序带来24格M改善，却仍存在3格额外COPY增加，因此没有全指标支配关系。

C04六个预选格全部退化，否定了对该固定构造的收益预期，却不足以否定整个零spill或联合通信构造方向。R9F三个候选绝对M改善但G下降，揭示的是接受政策与主要目标的冲突；它们未被完整策略接受，不能手工加入主表。这些负结果使方法选择的依据更清楚，也防止只展示成功图片导致过度归因。

\subsection{下界、可证差距与理论未覆盖部分}\label{ux4e0bux754cux53efux8bc1ux5deeux8dddux4e0eux7406ux8bbaux672aux8986ux76d6ux90e8ux5206}

当必要下界与方案上界来自同一执行域时，可用\(U/L-1\)给出相对最优差距的保守上界。P2已有冻结证书报告：在1～5核下，满足\(U\le1.05L\)的图数分别为81、62、32、33、23。这个统计意味着相应格可在该证书条件下认证距最优不超过5\%；它不是在500个互相独立样本上的统计推断。

5核的平均\(B/L\)约为6.108，明显高于实现的4.550。二者的差既可能来自构造尚有改进空间，也可能来自下界忽略真实通信和容量约束；不能把6.108写成已经证实可达的性能目标。附录给出的强化窗口与残余等待组合是数学推导，尚未生成对应的新全量证书，故不替换这里的冻结数值。

\begin{figure}[!htbp]
\centering
\includegraphics[width=\linewidth,height=.84\textheight,keepaspectratio]{figures/paper/fig-bound-gap.pdf}
\caption{同一优化域中的可行结果与必要界。$L\le\mathrm{OPT}\le U$，阴影是必要界给出的上限与当前结果之间的差距，不代表一定可实现的收益。}
\label{fig:bound-gap}
\end{figure}

\subsection{有效性边界}\label{ux6709ux6548ux6027ux8fb9ux754c}

全部100个官方图均参与了研发，本文结果属于给定公开评测集合上的完整证据，而非未见分布上的泛化测试。同一图的五个核数相互关联，不能把500格直接当作500个独立随机样本计算显著性。算法按输入结构而非case编号选策略能够减少显式查表，但不消除设计过程中接触公开图所带来的适配性。

本文没有将零散历史最优合并为主成绩，也没有把某次局部评分循环加速推断为完整求解器相同比例的加速。程序失败、超时与未运行应逐格保留；当前主批次完整成功，并不意味着所有未来输入都不会失败。最后，官方事件模拟是统一比较条件，其结果不等于已在真实NPU上验证的端到端收益。

\FloatBarrier

\section{模型评价与结论}\label{ux6a21ux578bux8bc4ux4ef7ux4e0eux7ed3ux8bba}

\subsection{方法的优势}\label{ux65b9ux6cd5ux7684ux4f18ux52bf}

本文以同一个可提交接口和执行语义层次连接三问。问题一围绕Task边界处理屏障与并行度；问题二利用同核驻留，把张量复用引入分核费用并与日历构造结合；问题三进一步安排归约次序和共享输入访问，在Cache语义下重新评价。三种方法共同采用结构守卫、有限候选、必要条件与完整评价，使每一项简化都有明确用途和验证边界。

数学分析直接服务于算法选择。秩函数解释为什么局部援助不制造等待环；超边触核费用解释为什么联合移动可能优于逐点迁移；树前沿交换证明给出无需枚举兄弟排列的次序；候选下界避免为不可能获胜的方案支付完整评价成本。这些结论与代码中实际检查的条件相连，而不是独立放置的理论装饰。

实验保留完整固定版本及逐格原件，区分方案质量与求解速度。真实JSON、子图映射和事件时间线使用同源标识，使机制解释能够追踪到具体执行。负例与代价同时报告，有助于说明为何不能仅用搬运量、命中率或局部负载衡量算法。

\subsection{局限与可推广条件}\label{ux5c40ux9650ux4e0eux53efux63a8ux5e7fux6761ux4ef6}

首先，完整执行仍然具有动态耦合。共享DDR、Cache填充/淘汰、FIFO和物理存储申请会相互影响，本文没有证明一个固定权重DAG足以精确表示全部行为，也未构造一般未来充分状态。局部编译复用、响应表示和完整方案去重应继续分别检查适用条件。

其次，结构识别限制了专用构造的适用范围。私有树、单输出、唯一生产者或无别名条件被破坏后，已有局部最优性不能直接延续；需要重新定义状态、证明或退回一般构造。容量、带宽、Core 数和官方核内策略改变时，也需重新校验费用模型、候选界和固定预算下的表现。

再次，当前主策略仍有可改进空间。R9F的比值门槛排除了绝对Makespan更好的局部候选，说明辅助指标应服务主要目标；P2中通信量与时间的取舍尚不能完全由简洁代理解释；P1中细粒度并行仍可能受大额边界延迟限制。对这些问题，更合理的后续工作是围绕具体反例修正构造和接受规则，再为新的固定算法建立完整证据，而不是事后拼接局部最佳数值。

最后，本文的数据来自给定的公开图与官方模拟，尚无未见分布、不同硬件条件和真实设备的效果验证。推广需要相应的新证据，不能从当前完整评测直接推出跨平台收益。

\subsection{结论}\label{ux7ed3ux8bba}

针对多核NPU的三种执行场景，本文形成了语义一致的结构化切分与调度方法。问题一通过屏障感知候选与有秩证书的支路援助组织并行，问题二以复用超图和空隙插入联合处理通信及计算，问题三将私有归约树的存储次序与官方Cache执行闭环结合。三个固定算法均完成100图×1～5核的完整评测，5核平均基线加速比分别为4.0553、4.5498和4.7576；P3同计划Cache加速比均值为1.0829。

结果支持本文的核心判断：图上的结构改动只有经过相应机器语义的解释，才能确定实际收益。有效的数学简化应保留当前决策所需的信息，并清楚声明不能推断的部分。由此形成的统一研究框架，将定义、局部证明、可执行构造、真实流水和完整实验连接为一条可检查的论证链。

\FloatBarrier
\begin{thebibliography}{9}
\bibitem{official} 第二十三届中国研究生数学建模竞赛. A题：通用神经网络处理器下的多核调度问题及官方附件[Z]. 2026.
\bibitem{graham} GRAHAM R L. Bounds for Certain Multiprocessing Anomalies[J]. Bell System Technical Journal, 1966, 45(9): 1563--1581. DOI: 10.1002/j.1538-7305.1966.tb01709.x.
\bibitem{hypergraph} ÇATALYÜREK Ü V, AYKANAT C. Hypergraph-Partitioning-Based Decomposition for Parallel Sparse-Matrix Vector Multiplication[J]. IEEE Transactions on Parallel and Distributed Systems, 1999, 10(7): 673--693. DOI: 10.1109/71.780863.
\bibitem{sethi} SETHI R, ULLMAN J D. The Generation of Optimal Code for Arithmetic Expressions[J]. Journal of the ACM, 1970, 17(4): 715--728. DOI: 10.1145/321607.321620.
\end{thebibliography}

\appendix
\section{补充数学推导}\label{ux8865ux5145ux6570ux5b66ux63a8ux5bfc}

\subsection{成功执行的必要计算投影}\label{ux6210ux529fux6267ux884cux7684ux5fc5ux8981ux8ba1ux7b97ux6295ux5f71}

在唯一原生产者、无原始逻辑身份混叠等保守条件下，只保留原计算之间直接存在的依赖，及同一张量原生产者到消费者的关系，不随意收缩经过原COPY的路径。任意官方成功执行保留全部原始计算及其时长，每核每Pipe仍为单槽。取真实轨迹中这些计算的原核及起止时刻，得到投影图的一个可行调度。

跨Task或跨核的因果不能只看一个被删除的逻辑边。P1通过前驱Task完成与激活门控保留因果；P2/P3通过源OUT、基础IN和500周期释放保留因果。接收侧后续reload还受Step2顺序及MTE2 FIFO制约，不可越过必要的基础IN。该投影论证说明保留的约束对每条成功轨迹成立；它不是把真实约束删掉后重新运行一个贪心调度器。

\subsection{非空窗口中的工作量与不可分原子}\label{ux975eux7a7aux7a97ux53e3ux4e2dux7684ux5de5ux4f5cux91cfux4e0eux4e0dux53efux5206ux539fux5b50}

在必要计算投影上，定义不含自身的最长前缀和后缀

\[h_v=\max_{u\to v}(h_u+d_u),\qquad t_v=\max_{v\to w}(d_w+t_w).\]

取非负\(a,b\)及某一Pipe的非空集合

\[S=\{v:p_v=p,\ h_v\ge a,\ t_v\ge b\}.\]

这些操作在任何成功执行中完整落在\([a,M-b]\)内。按时长降序记\(d_{(1)}\ge\cdots\ge d_{(n)}\)，令

\[A_K(S)=\max_{\substack{r\ge1\\m=(r-1)K+1\le n}}\sum_{j=m-r+1}^{m}d_{(j)}.\]

则有

\[M\ge a+b+\max\left\{\left\lceil\frac{\sum_{v\in S}d_v}{K}\right\rceil,A_K(S)\right\}.\]

\textbf{证明。} 窗口长度为\(x=M-a-b\)，同Pipe的\(K\)个槽总工作容量为\(Kx\)，给出工作量项。对前\(m=(r-1)K+1\)个操作，由鸽巢原理至少一个 Core分到\(r\)个；即使选取其中最短的\(r\)个，该核心也必须占用\(\sum_{j=m-r+1}^{m}d_{(j)}\)。逐\(r\)取最大得到原子项。两者是同一窗口的必要条件，取最大值而不是相加。最后只加一次头尾。\(\square\)

例如两个 Core上三个不可分操作时长为10、8、7，工作量给出13，简单均匀原子界可能只给14，而上述界给出15。这个例子只解释不可分性；是否能在真实机器达到15，还取决于其他依赖和资源。

计算head/tail为\(O(n+m)\)；给定窗口后，成员检查与时长排序分别为\(O(n)\)及\(O(|S|\log|S|)\)。这不意味着遍历所有双阈值窗口也具有同样复杂度。空集合不能贡献\(a+b\)。

\subsection{通用计算屏障与残余等待}\label{ux901aux7528ux8ba1ux7b97ux5c4fux969cux4e0eux6b8bux4f59ux7b49ux5f85}

取一条有序的通用计算屏障链；``通用''指这些计算与相关其他计算均由有向路径可比较，使屏障间的计算能够定位在互不重叠的实际时间块中。跨核必要等待记为\(\tau_q\)，P1为1000，P2/P3为500。该常数仅用于已有因果证明的跨核计算路径，不把可提前预读的输入事件强塞到计算块中。

对块内一个非空同Pipe窗口，头尾阈值为\(a,b\)，左、右是否存在屏障分别为\(\ell,r\in\{0,1\}\)。头尾路径可能已经覆盖部分跨核等待，尚需保留的等待仅为

\[c(a,b)=\ell(\tau_q-a)_++r(\tau_q-b)_+.\]

设块长为\(x\)，窗口长\(z=x-a-b\)，安全的容量上包络为

\[W(S)\le z+(K-1)(z-c)_+.\]

两端屏障同核时，一个 Core可利用全部\(z\)，其他核至多利用\((z-c)_+\)。两端异核时，令\(A=(\tau_q-a)_+\)、\(B=(\tau_q-b)_+\)，总容量至多

\[ (z-A)_++(z-B)_++(K-2)(z-A-B)_+.\]

由对\(z\le\min(A,B)\)、中间区间及\(z\ge A+B\)分段，可验证

\[(z-A)_++(z-B)_+\le z+(z-A-B)_+,\]

故仍受相同上包络约束。整数容量不等式\(W\le z+(K-1)(z-c)_+\)的最小可行长度为

\[\phi_{K,c}(W)=\min\left\{W,\left\lceil\frac{W+(K-1)c}{K}\right\rceil\right\}.\]

当\(W\le c\)时只用首核的\(z=W\)即可满足放松容量；当\(W>c\)时解第二段线性不等式并取整。结合B.2原子项，得到

\[x\ge a+b+\max\{\phi_{K,c(a,b)}(W(S)),A_K(S)\}.\]

各块与屏障区间互不重叠，故可将每块已证下界与屏障计算时间相加。仅仅从若干互相重叠窗口取得下界时不能相加，应取最大值。P1的同核100周期等待也不能任意加进来，因为一个计算屏障并不强迫方案在此处切开Task。

本式的要点是扣除头尾已覆盖的等待。把\(a+b\)、完整双向等待与工作量简单全部相加，会重复计费。该强化组合用于理论分析；本文主表仍引用既有冻结证书，没有声称已为新公式重新生成500格数值。

\subsection{未来充分性与当前评分表示}\label{ux672aux6765ux5145ux5206ux6027ux4e0eux5f53ux524dux8bc4ux5206ux8868ux793a}

第3章给出的三条件说明，动作是否可用、压缩状态怎样更新和当前观测怎样恢复，都必须由压缩状态决定。只满足第三项是当前评分表示；保存一个模拟器的运行现场是runtime checkpoint；缓存同一局部编译产物是编译复用。它们各自可以有效节省计算，但不自动满足关于所有未来动作词的等价要求。

固定观测也很重要。若完整结果能恢复计划标签与序列，要求``完整结果相同''可能让等价类极细，失去压缩收益；若只保留Makespan，当前等价类又可能在后续切分或移动后分裂。研究应从实际想省掉的工作和需要保留的观测出发，而不预设必须存在很小的全局状态空间。

\clearpage
\section{固定实现与复现信息}
三题固定求解器的提交分别为下列SHA。逐格图、计划、结果来源及完整SHA-256见配套CSV与来源收据。固定版本结果不能用不同算法的历史赢家拼接。
\begin{itemize}
\item P1：\texttt{\seqsplit{834d8c957538ee069c66aadac9509552a4cc69d7}}；在线评分至多10次E1。
\item P2：\texttt{\seqsplit{c66559a6f8a31ef7b4720e1f7c3c28d61f8dff3f}}；在线评分至多3次E2。
\item P3：\texttt{\seqsplit{311322b996c0948e8a6a9c7ec6ddfe6ae41fbee1}}；在线评分至多3次E0。
\end{itemize}
官方代码与配置身份分别为\texttt{\seqsplit{de11a83db8d7c47ed328b15a7df71d613a833b16cd23ee9fe877999578a1ace0}}和\texttt{\seqsplit{dcd10de54b23f8366428fb24e828812b1da9549e6eae4a3c3f38604fe5ae77b9}}。来源收据记录822个固定Git对象；本稿的数据导出仅读取已有原件，没有新增求解器或评价器调用。

\end{document}
